% !TEX root = ../main.tex
\section{Conclusiones}
\label{sec:conclusiones}

A partir del desarrollo formal, la implementación declarativa en MiniZinc y la evaluación experimental sobre el caso base y las diez instancias de planificación energética, se establecen las siguientes conclusiones técnicas justificadas por los resultados obtenidos:

\begin{enumerate}
    \item \textbf{Importancia crítica del acoplamiento dual en la función objetivo de un COP lineal entero:}
    En problemas de optimización donde la función objetivo maximiza la diferencia entre ingresos comerciales y costos de producción ($Z = \sum P_s vendido_{s,j} - \sum c_i prod_{i,j}$), la formulación primaria por sí sola presenta una relajación de intervalos débil, pues el propagador de la ecuación objetivo evalúa el máximo ingreso asumiendo simultáneamente producción nula ($\underline{prod}_{i,j} = 0$). La incorporación de la igualdad de flujo diario sin sobreproducción estéril ($R_2$), la dominancia por orden de mérito a través de la frontera de régimen normal ($R_4$) y las proyecciones lineales por escalón de costo marginal ($R_5$) acopló algebraicamente el costo mínimo de generación al margen neto por megavatio vendido $(P_s - c_q)$. Como se demostró en la Tabla~\ref{tab:ablacion_cop}, esta reformulación transformó un modelo que sufría \emph{timeout} ($>2.5\text{ s}$) incluso en $18$ variables ($M_0/M_1$) en un modelo que resuelve \texttt{Datos.dzn} determinísticamente en el nodo raíz ($1$ nodo, $0$ fallos, $0.09\text{ ms}$) y escala hasta instancias de $168$ variables con tres pares de plantas gemelas (\texttt{Instancia10.dzn}).

    \item \textbf{Generalidad estructural y exactitud de la restricción de régimen alto en todo el parque:}
    La corrección de la ventana deslizante de longitud exacta $R_i$ días ($k \in \{1, \dots, n - R_i + 1\}$ con $d \in \{0, \dots, R_i - 1\}$) aplicada uniformemente a todas las centrales eléctricas $i \in P$, complementada por la cota lineal acumulada por ventana $R_3$ \eqref{eq:r3_ventana_lineal}, permitió capturar con fidelidad el desgaste diferencial de turbinas nucleares, hidroeléctricas, térmicas y renovables. Los experimentos en \texttt{Instancia02.dzn} e \texttt{Instancia06.dzn} evidenciaron que el modelo planifica automáticamente esquemas de alternancia en contrafase entre centrales (operando un día a capacidad nominal $Cap_i$ en régimen alto y enfriando la turbina al día siguiente exactamente en su umbral normal $U_i$), y que cuando el costo marginal de la central de respaldo supera la tarifa de un cliente ($P_s < c_i$ en \texttt{Instancia03.dzn}, \texttt{Instancia04.dzn}, \texttt{Instancia06.dzn}, \texttt{Instancia08.dzn} e \texttt{Instancia10.dzn}), recorta de manera óptima la entrega de dicho cliente hacia su piso garantizado $G\%$ para evitar pérdidas operativas.

    \item \textbf{Eficacia del rompimiento de simetrías en parques con unidades modulares equivalentes:}
    El análisis de simetrías demostró que estás sólo existen dos o más centrales comparten idéntica tripla $(c_i, Cap_i, R_i)$ o cuando dos clientes comparten tarifa y perfil de demanda. La imposición de restricciones globales \texttt{lex\_greatereq} ($S_1$ y $S_2$) junto con el desempate canónico intra-día eliminó las particiones equivalentes de carga sin introducir sobrecosto en instancias heterogéneas, logrando en Gecode una reducción del $29.6\%$ de los nodos en \texttt{Instancia06.dzn} ($432 \to 304$) y del $33.2\%$ en \texttt{Instancia10.dzn} ($143\,390 \to 95\,808$ nodos y $9.93\times 10^6$ propagaciones ahorradas).

    \item \textbf{Superioridad del orden causal de decisión y la bisección descendente de intervalos:}
    El estudio comparativo de ocho estrategias de búsqueda (Tabla~\ref{tab:estrategias_cop}) refutó la idoneidad ingenua de la heurística clásica \texttt{first\_fail} + \texttt{indomain\_min} ($E_6$) para este COP: por una parte, la asimetría de escalas entre los dominios estrechos de $vendido_{s,j}$ y los dominios amplios de $prod_{i,j}$ hace que \texttt{first\_fail} ramifique sobre las ventas antes de resolver el cuello de botella inter-temporal de las plantas. Por otra parte, \texttt{indomain\_min} explora primero las asignaciones de peor ganancia ($\underline{Z}$), anulando la poda de \emph{Branch-and-Bound}. La estrategia adoptada $E_1$ (\texttt{input\_order} sobre $prod \mathbin{+\!+} vendido$ con \texttt{indomain\_reverse\_split}) resolvió la totalidad de las 11 instancias en ambos solvers, reduciendo en \texttt{Instancia08.dzn} los nodos explorados por Gecode en $17.4\times$ respecto a \texttt{indomain\_max} ($E_2$), en $77.3\times$ respecto a $E_6$ y en $99.0\times$ respecto a \texttt{indomain\_min} ($E_3$).

    \item \textbf{Validación de utilidad de técnica de generalidad:} De forma general la implementación del problema validó la gran utilidad y flexibilidad de solución que aporta la generalización de problemas de opmtimización por medio de la parametrización. Permitiendo en este caso construir un conjunto de restricciones que se adaptan a una familia de problemas de gran alta variedad como evidencian las pruebas realizadas. 
\end{enumerate}
